\section{Intersection numbers satisfy smoothing}
\label{sec:necessary}
In this section we give a detailed proof that intersection numbers satisfy the hypotheses of our theorem. The only non-trivial condition to check is smoothing. In \cite[4.1. Examples]{MGT21:Smoothings} we gave a topological proof of this fact for curves, and then conclude the result for any geodesic current by invoking density of weighted closed curves in the space of currents.
Here we give a different proof for an arbitrary current that works in the universal cover, since it will be useful to prove later that measured laminations satisfy max-smoothing.

\begin{proposition}
Given a geodesic current $\mu$, the curve functional $i(\mu, \cdot)$ satisfies the hypotheses of Theorem\ref{mainthm:intersections}.
\end{proposition}
\begin{proof}
\item (additivity and homogeneity) By definition (see \cite{Bonahon86:EndsHyperbolicManifolds}, intersection number is bilinear in both factors. In particular, additivity and homogeneity hold.
\item (stability) Observe that given a curve $C$ with hyperbolic lift $\gamma$, and $x \in \gamma$, the set of geodesics $G(x, \gamma^n x]$ can be partitioned into $\cup_{i=1}^{n-1} G(x, \gamma^{i} \cdot x]$.  Thus, by \cite[Proposition 3.14(ii)]{MZ19:PositivelyRatioed}, 
\[
i(\mu, \gamma^n) = \mu(G(x, \gamma^n x]) = \sum_{i=1}^{n-1} \mu(G(x, \gamma^{i} \cdot x]) = n\mu(G(x, \gamma x]) = n i(\mu, \gamma).
\]
where in the second inequality we used finite additivity of $\mu$, and in the third one, $\Gamma$-invariance.
\item (smoothing) Follows from Proposition \ref{prop:smoothing_nec}
\end{proof}

The following proof follows similar ideas as in the proof of \cite[Theorem A]{FF21:QF}.

\begin{proposition}
Suppose $a,b \in \Gamma$ are non-trivial elements with distinct fixed points on the boundary circle $\partial \Gamma$.
Let $\mu$ be any geodesic current.
Then,
\begin{enumerate}
\item If $b$ crosses $a$ to the right, then
\[
i(\mu, a) + i(\mu, b) \geq i(\mu, ab)
\]
and
\[
i(\mu, a) + i(\mu, b) \geq i(\mu, ab^{-1})
\]
\item If $a$ and $b$ are $R$-parallel, then
\[
 i(\mu, ab) \geq i(\mu, a) + i(\mu, b)
\]
and
\[
 i(\mu, ab) \geq i(\mu, ab^{-1})
\]
\item If $a$ and $b$ are $R$-anti-parallel, then
\[
 i(\mu, ab^{-1}) \geq i(\mu, a) + i(\mu, b)
\]
and
\[
 i(\mu, ab^{-1}) \geq i(\mu, ab)
\]
\end{enumerate}
\label{prop:smoothing_nec}
\end{proposition}

\begin{figure}[h!]
     \centering
     \begin{subfigure}[b]{0.45\textwidth}
\centering{
\resizebox{80mm}{!}{\svgfig{smoothnec1}}
    \caption{$b$ right crosses $a$}\label{fig:smoothing1}}
     \end{subfigure}
     \hfill
     \begin{subfigure}[b]{0.45\textwidth}
\centering{
\resizebox{80mm}{!}{\svgfig{smoothnec2}}
    \caption{$a$ and $b$ are $R$-parallel}\label{fig:smoothing2}}
     \end{subfigure}
\end{figure}

\DMGtodo[inline]{probably want to add the MZ result in the text}
\begin{proof}
\begin{enumerate}
\item We have a situation like in Figure \ref{fig:smoothing1}.
Let $p$ be the intersection point between the hyperbolic axis $A_a$ of $a$ and $A_b$ of $b$.
Since $p \in A_a$, we have $x=b^{-1} \cdot p \in A_b$.
We have $p$ and $y= a \cdot p$ are in $A_a$. Thus
\begin{align*}
i(\mu, ab) \leq & \mu(G(x, ab x]) \leq \mu(G(x, b x]) + \mu(G(bx, ab x]) = \\
&  = \mu(G(b^{-1} p, p]) + \mu(G(p, ap]) = i(\mu, b) + i(\mu, a)
\end{align*}
In the first inequality we used \cite[Proposition 3.14(ii)]{MZ19:PositivelyRatioed},, in the second, we used that 
$G(x, b x] \subset G(x, b x] \cup G(bx, ab x]$, and in the last equality we used the fact that $p$ is in the
axis of $a$ and of $b$, and \cite[Proposition 3.14(i)]{MZ19:PositivelyRatioed},
A similar argument with $ab^{-1}$ gives the other inequality.
\item We have a situation like in Figure \ref{fig:smoothing2} (note that the axes of $aB$ and $Ba$ can be in another relative position, but this will not matter).
The first inequality can be proven as follows.
Let $p$ be the intersection between $A_{ab}$ and $b \cdot A_{ab} = A_{ba}$.
Since $p \in b \cdot A_{ab}$, we have that $x = b^{-1} \cdot p$ is in $A_{ab}$ and thus $ab \cdot x \in A_{ab}$, also.
We thus have
\begin{align*}
i(\mu, ab) &= \mu(G(x, ab x]) =  \mu(G(x, p]) +  \mu(G(p, ab \cdot x]) = \\
 & =  \mu(G(b^{-1} p, p]) +  \mu(G(p, a p]) \geq i(\mu, a) + i(\mu, b).
\end{align*}
In the first equality we use \cite[Proposition 3.14(i)]{MZ19:PositivelyRatioed}, in the second equality we used the fact that $G(x, ab x] = G(x, p] \cup G(p, ab \cdot x]$
gives a partition of $G(x, ab x]$, and $\mu$ is finitely additive. The last inequality uses again \cite[Proposition 3.14(ii)]{MZ19:PositivelyRatioed},.
For the second inequality, let $y = a^{-1} \cdot p$, which is in $A_{ba}$, since $p \in a \cdot A_{ba}=A_{ab}$.
Thus, similarly as before, we get
\begin{align*}
i(\mu, ba) &= \mu(G(y, ba y]) =  \mu(G(y, p]) +  \mu(G(p, x]) = \\
 & =  \mu(G(a^{-1} \cdot p, p]) +  \mu(G(p, b^{-1} p]) \geq \mu(G(a^{-1} \cdot p, b^{-1} p]) \geq \\
i(\mu, b^{-1} a).
\end{align*}
The only inequality that is different from the arguments above is the second. We partitioned
$G(y, ba y]$ as $G(y, p] \cup G(p, ba y]$, and then applied $b^{-1}$ to $G(p, ba y]$ to obtain $G(p,x]$.
Since $\mu$ is $\Gamma$-invariant, the third equality then holds. The fourth inequality is because
$G(a^{-1} \cdot p, b^{-1} p] \subset G(a^{-1} \cdot p, p] \cup G(p, b^{-1} p]$. The last inequality follows by 
\cite[Proposition 3.14(i)]{MZ19:PositivelyRatioed}. This concludes the proof of the second inequality.
\item This case is the same as the second case, by changing $b$ to $b^{-1}$.
\end{enumerate}
\end{proof}